% !TEX root = ../main.tex

\chapter{绪论}
\chaptereng{Introduction}


\section{课题来源、选题依据和立论背景}

火焰温度的准确测量与标定是实现燃烧过程精细化控制的关键前提，在燃气轮机燃烧、钢铁冶炼、石英制造等关键工业场景中，火焰温度的精准调控直接关系到能源消耗、污染物排放与产品质量。工业生产对火焰标定提出了现场适用性、宽温度覆盖范围、高稳定性与可溯源性等核心要求，但传统火焰温度测量方法普遍存在局限性：接触式测量（如热电偶）易受燃烧环境腐蚀、辐射散热影响，无法复现工业燃烧的动态环境；非接触式光学测量（如激光诊断）虽避免了接触干扰，却缺乏统一的现场标定标准，导致测量结果难以溯源。这些问题严重制约了不同测量系统间结果的可比性与可靠性。

国际计量领域的技术革新为火焰标定突破提供了重要支撑。国际温标（ITS-90）的推广明确了温度测量的溯源性要求，欧洲计量创新与研究计划（EMPIR）资助的EMPRESS、EMPRESS 2等重点项目，将标准火焰标定技术推向实用化\cite{fateev2019validation}。

\begin{figure}[H]
	\centering
	
\includegraphics[width=0.4\textwidth]{image1.png}
	\caption{欧洲EMPRESS 2项目示意图}
	\label{fig:empress}
\end{figure}

标准火焰作为可溯源、高稳定性的标定源，是连接基础计量标准与实际应用的关键纽带。标准火焰用于标定的核心价值在于：为燃烧领域的温度测量提供可追溯的标定基准，解决传统测温技术没有标准可以依照的痛点；标准火焰作为一个桥梁可以将传统的热电偶测温与激光测量技术相连接，用标准火焰去标定激光测量的准确性；通过模拟真实燃烧环境的温度场与物种分布，为各类光学诊断技术提供精准的性能验证与误差校准，是优化测量精度的核心工具；通过提升温度测量精度，间接实现工业过程的精准控制，最终带来能耗降低、废物减少的经济效益与环境效益。

非预混平面绝热火焰炉的技术优势使其成为火焰标定的理想载体：其一，燃烧稳定性强，通过优化流体力学设计与氮气共流技术，可有效抑制外界气流干扰，确保温度场空间均匀性与长期重复性；其二，温度可溯源，借助激光瑞利散射等技术对炉内火焰温度进行精确表征，能够实现2050K-2250K温度范围的ITS-90溯源，测量不确定度可控制在0.5\%以内。以非预混平面绝热火焰炉作为火焰标定源，可构建多温度点的均匀温度场，为火焰标定提供了标准化的参考环境\cite{fateev2019validation}。

尽管火焰标定技术已取得显著进展，但仍存在待突破的理论与技术问题。其中高温极端环境下的温度覆盖范围有待拓展，现有以非预混平面绝热火焰炉作为标定源的温度上限大约为2300K，当进一步升高温度时就会导致温度场的不均匀性增大，且内部结构容易出现烧蚀的问题；现有的火焰光谱仿真大都是用国际温标的数据，而国际温标中的数据基本都是通过电加热炉加热到预定温度后用热电偶进行测量，因此在光谱展宽等方面与真实火焰会有一定的差异，实验数据与之前数据库的仿真存在着一定的差异\cite{liu2018two}。因此，开展基于非预混平面绝热火焰炉的火焰标定研究，建立覆盖更广温度范围、适应复杂环境、具备溯源性的火焰标定体系，能为我国高端制造业的燃烧过程精准控制提供核心技术支撑，对降低能源消耗、减少污染物排放、提升产品质量具有重要的理论价值\cite{qinghui2023experimental}。

\begin{figure}[H]
	\centering
	\begin{subfigure}{0.45\textwidth}
		\centering
		
\includegraphics[width=\linewidth]{image2.png}
		\label{fig:erosion}
	\end{subfigure}
	\hfill
	\begin{subfigure}{0.45\textwidth}
		\centering
		\includegraphics[width=\linewidth]{image3.png}
		\label{fig:spectrum_deviation}
	\end{subfigure}
	\caption{高温烧蚀照片与实验数据与仿真拟合的差异\cite{liu2018two}}
	\label{fig:erosion_comparison}
\end{figure}

在全球双碳目标推进与能源结构转型的大背景下，火焰标定的工业应用价值愈发凸显，尤其在燃气轮机高效运行场景中，其作为连接计量标准与工业应用的核心环节，能够建立可溯源、高稳定性的燃烧状态基准。基于标准火焰标定的光学诊断技术，可实现燃烧室温度场的精准监测，破解燃气轮机燃烧控制精度不足的难题，进而延长涡轮叶片寿命、提升燃气轮机热效率；同时，燃气轮机常用的贫燃预混燃烧技术需精准控制当量比以降低NO$_x$排放，而当量比偏差易引发燃烧不稳定，标准火焰恰好可用于验证燃气轮机低NO$_x$燃烧器性能，通过对比红外成像与标准火焰温度数据，降低当量比控制误差，提升NO$_x$排放达标率\cite{liu2022investigation}。

除燃气轮机领域外，火焰标定还能显著提升氢能燃烧的稳定性与安全性。氢能燃烧具有燃烧速度快、温度高、点火能量低、易发生回火等特性，实际应用中常采用与天然气掺混的方式使用，但掺混过程中燃料组分波动会直接导致火焰当量比偏移，引发燃烧不稳定问题\cite{mo2024氢气}。火焰标定技术可精准模拟不同氢能掺混比例的燃烧场景，通过标定温度与物种浓度，为氢能燃烧系统提供统一标准。

\section{课题研究目的、理论意义和实际应用价值}

\subsection{课题研究目的}

针对当前非预混平面绝热火焰炉作为标定源仍然缺少对其火焰结构的详细测量和不能覆盖高温极端环境的问题：

第一，重新设计非预混平面绝热火焰炉结构，添加水冷机构，更改喷嘴的排布，解决其不能在高温极端环境下稳定长久工作的问题，使非预混平面绝热火焰炉燃烧氢氧火焰温度更接近绝热温度，拓展高温标定范围。

第二，研究微扩散火焰结构，建立更精确的火焰温度场浓度场，对以非预混平面绝热火焰炉作为标定源的火焰标定提供更为详细的参考，同时为后续标定提供一个可溯源的标准化标定条件，也为后续工业场景标定应用夯实基础。

\subsection{理论意义}

使非预混平面绝热火焰炉作为标定源可以用于标定高温环境，在燃烧氢氧火焰时温度可以达到3000K，解决当下以非预混平面绝热火焰炉作为标定源所进行标定的研究没有覆盖高温环境的问题，和缺少对火焰结构的详细测量的问题，完善燃烧温度测量的溯源体系，为后续标定提供一个可溯源的标准化标定条件，为燃烧过程的精准控制提供技术支撑，推动燃烧计量领域的理论发展与技术革新。

\subsection{实际应用价值}

提高氢能燃烧的稳定性与安全性，通过标定明确氢气燃烧的临界条件、贫燃熄火极限等关键参数，为氢能燃烧设备的结构设计提供数据支撑，降低安全风险。在燃气轮机领域，可以精准控制污染物排放，通过标准火焰验证测温技术，确保燃气轮机燃烧温度控制在1800K以下的NO$_x$抑制区间，减少热损失，提高燃烧室热效率，延长叶片使用寿命。

\section{论文结构安排}

本文共分六章。第一章为引言，阐述研究背景与目标。第二章为文献综述，系统梳理平面火焰炉、TDLAS测量与火焰标定技术的国内外研究现状。第三章详述试验台与硬件系统，包括新型炉体设计、TDLAS光学链路与探测系统。第四章介绍研究方法与数据处理算法。第五章呈现实验结果与讨论。第六章为结论与展望。
